\begin{thebibliography}{99}
\raggedright

\bibitem{AhagCzyzVirtanen2026}
P.~\AA hag, R.~Czy{\.{z}}, and J.~Virtanen,
\emph{The complete Crouzeix conjecture in dimension three and the
Clou\^atre--Ostermann--Ransford conjecture},
\href{https://arxiv.org/abs/2608.27346v4}{arXiv:2608.27346v4}, September 21, 2026.

\bibitem{Ahlfors1979}
L.~V. Ahlfors, \emph{Complex Analysis}, third ed., McGraw--Hill, New York, 1979.

\bibitem{Arveson1972}
W. Arveson, Subalgebras of $C^*$-algebras II,
\emph{Acta Math.} \textbf{128} (1972), 271--308.

\bibitem{BadeaCrouzeixDelyon2006}
C. Badea, M. Crouzeix, and B. Delyon, Convex domains and $K$-spectral sets,
\emph{Math. Z.} \textbf{252} (2006), 345--365.

\bibitem{Beckermann2005}
B. Beckermann, Image num\'erique, GMRES et polyn\^omes de Faber,
\emph{C. R. Math.} \textbf{340} (2005), 855--860.

\bibitem{BeckermannCrouzeix2014}
B. Beckermann and M. Crouzeix, Faber polynomials of matrices for non-convex sets,
\emph{Jaen J. Approx.} \textbf{6} (2014), 219--231.

\bibitem{BeckermannStylianopoulos2018}
B. Beckermann and N. Stylianopoulos, Bergman Orthogonal Polynomials and the
Grunsky Matrix, \emph{Constr. Approx.} \textbf{47} (2018), 211--235.

\bibitem{Brodskii1978}
M. S. Brodski\u\i, Unitary operator colligations and their characteristic
functions, \emph{Russian Math. Surveys} \textbf{33}:4 (1978), 159--191.

\bibitem{Crouzeix2004}
M.~Crouzeix, Bounds for analytical functions of matrices,
\emph{Integral Equations Operator Theory} \textbf{48} (2004), 461--477.
\href{https://doi.org/10.1007/s00020-002-1188-6}{doi:10.1007/s00020-002-1188-6}.

\bibitem{Crouzeix2007}
M.~Crouzeix, Numerical range and functional calculus in Hilbert space,
\emph{J. Funct. Anal.} \textbf{244} (2007), 668--690.
\href{https://doi.org/10.1016/j.jfa.2006.10.013}{doi:10.1016/j.jfa.2006.10.013}.

\bibitem{CrouzeixPalencia2017}
M. Crouzeix and C. Palencia, The numerical range is a
$(1+\sqrt2)$-spectral set,
\emph{SIAM J. Matrix Anal. Appl.} \textbf{38} (2017), 649--655.

\bibitem{DelyonDelyon1999}
B. Delyon and F. Delyon, Generalization of von Neumann's spectral sets
and integral representation of operators,
\emph{Bull. Soc. Math. France} \textbf{127} (1999), 25--41.

\bibitem{Hausdorff1919}
F. Hausdorff, Der Wertvorrat einer Bilinearform,
\emph{Math. Z.} \textbf{3} (1919), 314--316.

\bibitem{Jin2026}
S.~Jin, \emph{The numerical range is a $2$-spectral set},
\href{https://www.preprints.org/manuscript/202607.1919/v4}{Preprints.org 202607.1919, version 4},
August 7, 2026.

\bibitem{KatsnelsonKheifetsYuditskii1997}
V. E. Katsnelson, A. Ya. Kheifets, and P. M. Yuditskii,
An abstract interpolation problem and the extension theory of isometric operators,
in \emph{Topics in Interpolation Theory}, OTAA \textbf{95},
Birkh\"auser, 1997, 283--298.
Preprint version: \href{https://arxiv.org/abs/0706.1896v1}{arXiv:0706.1896v1}.

\bibitem{LoristSchwenninger2026}
E.~Lorist and F.~L. Schwenninger,
\emph{A solution to Crouzeix's conjecture},
\href{https://arxiv.org/abs/2608.03841v2}{arXiv:2608.03841v2}, August 17, 2026.

\bibitem{Luo2026}
Q.~Luo, \emph{Two boundary-certificate proofs of the Crouzeix theorem:
Fiberwise Fourier recurrence and defect-Gram telescoping},
\href{https://www.preprints.org/manuscript/202608.1661/v2}{Preprints.org 202608.1661, version 2},
August 28, 2026.

\bibitem{Direct}
OpenAI, \emph{A direct proof of the complete Crouzeix inequality},
OpenAI Math Release preprint
\href{https://github.com/openai/math/blob/main/preprints/A-direct-proof-of-the-complete-Crouzeix-inequality-September-26-2026/paper.pdf}{OAI:A-direct-proof-of-the-complete-Crouzeix-inequality-September-26-2026},
2026.

\bibitem{Paulsen1984Similarity}
V. I. Paulsen, Completely bounded homomorphisms of operator algebras,
\emph{Proc. Amer. Math. Soc.} \textbf{92} (1984), 225--228.

\bibitem{Pommerenke1992}
C. Pommerenke, \emph{Boundary Behaviour of Conformal Maps},
Springer-Verlag, 1992.

\bibitem{RansfordSchwenninger2018}
T.~Ransford and F.~L. Schwenninger,
Remarks on the Crouzeix--Palencia proof that the numerical range is a
$(1+\sqrt{2})$-spectral set,
\emph{SIAM J. Matrix Anal. Appl.} \textbf{39} (2018), 342--345.
\href{https://doi.org/10.1137/17M1143757}{doi:10.1137/17M1143757}.

\bibitem{Rockafellar1970}
R. T. Rockafellar, \emph{Convex Analysis}, Princeton University Press, 1970.

\bibitem{Runge1885}
C.~Runge, Zur Theorie der eindeutigen analytischen Functionen,
\emph{Acta Math.} \textbf{6} (1885), 229--244.
\href{https://doi.org/10.1007/BF02400416}{doi:10.1007/BF02400416}.

\bibitem{Stein1952}
P. Stein, Some General Theorems on Iterants,
\emph{J. Res. Nat. Bur. Stand.} \textbf{48} (1952), 82--83.

\bibitem{SzNagy1953}
B. Sz.-Nagy, Sur les contractions de l'espace de Hilbert,
\emph{Acta Sci. Math. (Szeged)} \textbf{15} (1953), 87--92.

\bibitem{Toeplitz1918}
O. Toeplitz, Das algebraische Analogon zu einem Satze von Fej\'er,
\emph{Math. Z.} \textbf{2} (1918), 187--197.

\end{thebibliography}
